% GENERATED FILE. Edit canonical lessons, then run npm run generate:books.
% canonical-source-hash: fnv1a64:85fe1115ed817461
% canonical-lessons: SA-C481-andhah, SA-C481-badhirah, SA-C481-mukah, SA-C481-panguh, SA-C481-kubjah

\chapter{Blind, Deaf and Lame}
\label{ch:sa-blind-deaf-and-lame}
\hlchaptermodality{481}

\begin{chapteropening}
\textbf{By the end of this chapter:} \emph{I can talk about someone who cannot see, hear or walk well.}

Complete the last lesson of chapter 481: I can talk about someone who cannot see, hear or walk well.
\end{chapteropening}

\section[\texorpdfstring{\sk{अन्धः} (andhaḥ)}{अन्धः (andhaḥ)}]{\sk{अन्धः} (andhaḥ) — blind}

\label{lesson:SA-C481-andhah}



\begin{quote}
Before the new one: say the Sanskrit for a broken bone, then the Sanskrit for oozes.
\end{quote}

\subsection*{You'll want to know: \sk{अन्धः}}
\textbf{\sk{अन्धः}} — \emph{andhaḥ} — \textquotedblleft{}blind\textquotedblright{}.

\begin{quote}
\sk{अन्धः} \sk{याचकः} \sk{मार्गे} \sk{गायति।}
\end{quote}

\begin{quote}
\emph{andhaḥ yācakaḥ mārge gāyati.}
\end{quote}

\begin{quote}
A blind beggar sings by the road.
\end{quote}

\begin{cousinweb}
\textbf{\sk{अन्धः}} (\emph{andhaḥ}) MW gives \textquotedblleft{}blind\textquotedblright{} and \textquotedblleft{}dark\textquotedblright{}; MW gives no derivation, so its origin is uncertain (MW).
\end{cousinweb}

\begin{grammarlens}[title={What you've built}]
One more word for describing what you lived and felt.
\end{grammarlens}

\subsection*{Your turn}
\begin{itemize}
\raggedright
  \item \emph{Say it:} \emph{andhaḥ}
  \item \emph{Say it:} \emph{andhaḥ}, once more
  \item \emph{Say it:} \emph{andhaḥ yācakaḥ mārge gāyati.} — A blind beggar sings by the road.
  \item \emph{Recall:} say the Sanskrit for a broken bone, then the Sanskrit for oozes, then say \emph{andhaḥ} again
\end{itemize}

\begin{tcolorbox}[breakable,colback=teal!4,colframe=teal!35!black,title={Before you move on}]
What does \sk{अन्धः} mean? (\textquotedblleft{}Blind\textquotedblright{}.) Say it once more.
\end{tcolorbox}
\section[\texorpdfstring{\sk{बधिरः} (badhiraḥ)}{बधिरः (badhiraḥ)}]{\sk{बधिरः} (badhiraḥ) — deaf}

\label{lesson:SA-C481-badhirah}



\begin{quote}
Before the new one: say the Sanskrit for oozes, then the Sanskrit for blind.
\end{quote}

\subsection*{You'll want to know: \sk{बधिरः}}
\textbf{\sk{बधिरः}} — \emph{badhiraḥ} — \textquotedblleft{}deaf\textquotedblright{}.

\begin{quote}
\sk{मम} \sk{बधिरः} \sk{पितामहः} \sk{उच्चैः} \sk{वदति।}
\end{quote}

\begin{quote}
\emph{mama badhiraḥ pitāmahaḥ uccaiḥ vadati.}
\end{quote}

\begin{quote}
My deaf grandfather talks loudly.
\end{quote}

\begin{cousinweb}
\textbf{\sk{बधिरः}} (\emph{badhiraḥ}) MW gives \textquotedblleft{}deaf\textquotedblright{} (sometimes written \emph{vadhira}); MW gives no derivation, so its origin is uncertain (MW).
\end{cousinweb}

\begin{grammarlens}[title={What you've built}]
One more word for describing what you lived and felt.
\end{grammarlens}

\subsection*{Your turn}
\begin{itemize}
\raggedright
  \item \emph{Say it:} \emph{badhiraḥ}
  \item \emph{Say it:} \emph{badhiraḥ}, once more
  \item \emph{Say it:} \emph{mama badhiraḥ pitāmahaḥ uccaiḥ vadati.} — My deaf grandfather talks loudly.
  \item \emph{Recall:} say the Sanskrit for oozes, then the Sanskrit for blind, then say \emph{badhiraḥ} again
\end{itemize}

\begin{tcolorbox}[breakable,colback=teal!4,colframe=teal!35!black,title={Before you move on}]
What does \sk{बधिरः} mean? (\textquotedblleft{}Deaf\textquotedblright{}.) Say it once more.
\end{tcolorbox}
\section[\texorpdfstring{\sk{मूकः} (mūkaḥ)}{मूकः (mūkaḥ)}]{\sk{मूकः} (mūkaḥ) — mute}

\label{lesson:SA-C481-mukah}



\begin{quote}
Before the new one: say the Sanskrit for blind, then the Sanskrit for deaf.
\end{quote}

\subsection*{You'll want to know: \sk{मूकः}}
\textbf{\sk{मूकः}} — \emph{mūkaḥ} — \textquotedblleft{}mute\textquotedblright{}.

\begin{quote}
\sk{मूकः} \sk{बालः} \sk{न} \sk{वदति}, \sk{किन्तु} \sk{सर्वं} \sk{जानाति।}
\end{quote}

\begin{quote}
\emph{mūkaḥ bālaḥ na vadati, kintu sarvaṃ jānāti.}
\end{quote}

\begin{quote}
The mute boy does not speak, but he understands everything.
\end{quote}

\begin{cousinweb}
\textbf{\sk{मूकः}} (\emph{mūkaḥ}) MW explains it as \textquotedblleft{}tied or bound\textquotedblright{}, that is, tongue-tied, and gives \textquotedblleft{}dumb, speechless, mute\textquotedblright{}; it names no root, so its origin is uncertain (MW).
\end{cousinweb}

\begin{grammarlens}[title={What you've built}]
One more word for describing what you lived and felt.
\end{grammarlens}

\subsection*{Your turn}
\begin{itemize}
\raggedright
  \item \emph{Say it:} \emph{mūkaḥ}
  \item \emph{Say it:} \emph{mūkaḥ}, once more
  \item \emph{Say it:} \emph{mūkaḥ bālaḥ na vadati, kintu sarvaṃ jānāti.} — The mute boy does not speak, but he understands everything.
  \item \emph{Recall:} say the Sanskrit for blind, then the Sanskrit for deaf, then say \emph{mūkaḥ} again
\end{itemize}

\begin{tcolorbox}[breakable,colback=teal!4,colframe=teal!35!black,title={Before you move on}]
What does \sk{मूकः} mean? (\textquotedblleft{}Mute\textquotedblright{}.) Say it once more.
\end{tcolorbox}
\section[\texorpdfstring{\sk{पंगुः} (paṅguḥ)}{पंगुः (paṅguḥ)}]{\sk{पंगुः} (paṅguḥ) — lame}

\label{lesson:SA-C481-panguh}



\begin{quote}
Before the new one: say the Sanskrit for deaf, then the Sanskrit for mute.
\end{quote}

\subsection*{You'll want to know: \sk{पंगुः}}
\textbf{\sk{पंगुः}} — \emph{paṅguḥ} — \textquotedblleft{}lame\textquotedblright{}.

\begin{quote}
\sk{पंगुः} \sk{वृद्धः} \sk{शनैः} \sk{चलति।}
\end{quote}

\begin{quote}
\emph{paṅguḥ vṛddhaḥ śanaiḥ calati.}
\end{quote}

\begin{quote}
The lame old man walks slowly.
\end{quote}

\begin{cousinweb}
\textbf{\sk{पंगुः}} (\emph{paṅguḥ}) MW asks, with a question mark, whether it comes from a root √paj, so its origin is uncertain; it gives \textquotedblleft{}lame, halt, crippled in the legs\textquotedblright{} (MW).
\end{cousinweb}

\begin{grammarlens}[title={What you've built}]
One more word for describing what you lived and felt.
\end{grammarlens}

\subsection*{Your turn}
\begin{itemize}
\raggedright
  \item \emph{Say it:} \emph{paṅguḥ}
  \item \emph{Say it:} \emph{paṅguḥ}, once more
  \item \emph{Say it:} \emph{paṅguḥ vṛddhaḥ śanaiḥ calati.} — The lame old man walks slowly.
  \item \emph{Recall:} say the Sanskrit for deaf, then the Sanskrit for mute, then say \emph{paṅguḥ} again
\end{itemize}

\begin{tcolorbox}[breakable,colback=teal!4,colframe=teal!35!black,title={Before you move on}]
What does \sk{पंगुः} mean? (\textquotedblleft{}Lame\textquotedblright{}.) Say it once more.
\end{tcolorbox}
\section[\texorpdfstring{\sk{कुब्जः} (kubjaḥ)}{कुब्जः (kubjaḥ)}]{\sk{कुब्जः} (kubjaḥ) — hunchbacked}

\label{lesson:SA-C481-kubjah}



\begin{quote}
Before the new one: say the Sanskrit for mute, then the Sanskrit for lame.
\end{quote}

\subsection*{You'll want to know: \sk{कुब्जः}}
\textbf{\sk{कुब्जः}} — \emph{kubjaḥ} — \textquotedblleft{}hunchbacked\textquotedblright{}.

\begin{quote}
\sk{कुब्जः} \sk{वृद्धः} \sk{यष्ट्या} \sk{चलति।}
\end{quote}

\begin{quote}
\emph{kubjaḥ vṛddhaḥ yaṣṭyā calati.}
\end{quote}

\begin{quote}
The hunchbacked old man walks with a staff.
\end{quote}

\begin{cousinweb}
\textbf{\sk{कुब्जः}} (\emph{kubjaḥ}) MW gives \textquotedblleft{}hump-backed, crooked\textquotedblright{}, from the Vājasaneyi Saṃhitā on; MW gives no derivation, so its origin is uncertain (MW).
\end{cousinweb}

\begin{grammarlens}[title={What you've built}]
One more word for describing what you lived and felt.
\end{grammarlens}

\subsection*{Your turn}
\begin{itemize}
\raggedright
  \item \emph{Say it:} \emph{kubjaḥ}
  \item \emph{Say it:} \emph{kubjaḥ}, once more
  \item \emph{Say it:} \emph{kubjaḥ vṛddhaḥ yaṣṭyā calati.} — The hunchbacked old man walks with a staff.
  \item \emph{Recall:} say the Sanskrit for mute, then the Sanskrit for lame, then say \emph{kubjaḥ} again
\end{itemize}

\begin{tcolorbox}[breakable,colback=teal!4,colframe=teal!35!black,title={Before you move on}]
What does \sk{कुब्जः} mean? (\textquotedblleft{}Hunchbacked\textquotedblright{}.) Say it once more.
\end{tcolorbox}
